\documentclass[11pt,a4paper]{amsart}
\usepackage[a4paper,landscape,margin=20mm]{geometry}
\usepackage{amsmath,amssymb,amsthm,graphicx,xcolor}
\usepackage[no-math]{fontspec}
\setmainfont{lmroman10-regular.otf}[BoldFont=lmroman10-bold.otf,ItalicFont=lmroman10-italic.otf,BoldItalicFont=lmroman10-bolditalic.otf,SmallCapsFont=lmromancaps10-regular.otf]
\usepackage[inline,shortlabels]{enumitem}
%% read in additional TeX-packages or personal macros here:
\DeclareMathOperator{\Reach}{Reach}
\newcommand{\step}[1]{\mathrel{\xrightarrow{#1}}}
\newcommand{\reach}{\step{*}}
\newcommand{\virtual}[1]{\mathrel{\xdashrightarrow{#1}}\hspace{-0.3em}}
\newcommand{\vreach}{\virtual{*}}
\newcommand{\vreachtitle}{\xdashrightarrow{*}}

\newcommand{\dist}[3]{(#1,#2)\text{-dist}(#3)}
\newcommand{\adj}{\text{adj}}
\newcommand{\sol}{\textsc{sol}}

\newcommand{\Cover}{\textsc{Cover}\xspace}
\newcommand{\SCover}{\textsc{SCover}\xspace}
\newcommand{\PPSCover}{\textsc{PPSCover}\xspace}

\newcommand{\direct}[1]{\stackrel{#1}{\Rightarrow}}

\newcommand{\sleq}{\preceq}
\newcommand{\minsleq}[1]{\min{}_\sleq(#1)}

\newcommand{\sign}{\textsc{sign}}
\newcommand{\esign}{\textsc{esn}}

\newcommand{\dif}{\textsc{dif}}
\newcommand{\eqstr}{\equiv_{\textsc{str}}}
\newcommand{\bvar}[1]{\mathbf{#1}}
\newcommand{\D}{{}_{\delta}}

\newcommand{\nnorm}[1]{||#1||}
\newcommand{\vr}{{}_\textnormal{vr}}
\newcommand{\transp}{\mathsf{T}}
\newcommand{\bvr}{{}_\textbf{vr}}
\newcommand{\mr}{{}_\textnormal{mr}}
\newcommand{\bmr}{{}_\textbf{mr}}
\newcommand{\upC}[1]{\mathop{\uparrow}#1}
\newcommand{\dead}{{}_\textnormal{dead}}
\newcommand{\bdead}{{}_\textbf{dead}}
\newcommand{\suf}{{}_\textnormal{suf}}
\newcommand{\bsuf}{{}_\textbf{suf}}
\newcommand{\shift}{{}_\textnormal{shift}}
\newcommand{\bshift}{{}_\textbf{shift}}
\newcommand{\close}{\textnormal{close}}
\newcommand{\qd}{{}_\textnormal{qd}}
\newcommand{\bqd}{{}_\textbf{qd}}

\newcommand{\ASC}{A_{\textsc{sc}}}

\newcommand{\length}[1]{|#1|}
\newcommand{\tu}[2]{\langle #1,#2\rangle}
\newcommand{\lara}[1]{{\langle #1\rangle}}
\newcommand{\ins}{\textsc{ins}}

\newcommand{\iotrans}[4]{#1:#2 \xrightarrow{#3} #4}
\newcommand{\pptrans}[3]{#1:(#2) \rightarrow (#3)}
\newcommand{\pptransone}[3]{#1: #2 \rightarrow #3}

\newcommand{\ccmp}[1]{\mathbb{CC}_{#1}}
\newcommand{\mone}{\mathbf{1}}
\newcommand{\preset}[1]{{^\bullet}#1}
\newcommand{\postset}[1]{#1^\bullet}
\newcommand{\premset}[1]{{^\bullet}#1}
\newcommand{\postmset}[1]{#1^\bullet}
\newcommand{\rst}[2]{#1{_{| #2}}}
\newcommand{\rset}[2]{\mathcal{R}(#1,#2)}
\newcommand{\SCC}{\mathcal{SCC}}
\newcommand{\ininter}[3]{#1 \in [#2,#3]}

\newcommand{\tr}{\textsc{tr}}
\newcommand{\dl}{_{\textsc{dl}}}
\newcommand{\bal}[2]{\textsc{bal}_{#2}{(#1)}}

\newcommand{\NP}{\textsc{NP}\xspace}
\newcommand{\coNP}{\textsc{co-NP}\xspace}
\newcommand{\PTIME}{\mathbf{P}}
\newcommand{\PSPACE}{\textsc{PSPACE}\xspace}
\newcommand{\NPSPACE}{\textsc{NPSPACE}\xspace}
\newcommand{\EXPTIME}{\mathbf{EXPTIME}}

\newcommand{\twoEXPTIME}{\textsc{2-ExpTime}}

\newcommand{\EXPSPACE}{EXPSPACE}
\newcommand{\ALTEXPSPACE}{\textsc{AltExpSpace}}
\newcommand{\TOWER}{\textsc{TOWER}}
\newcommand{\sizeof}[1]{|#1|}
\newcommand{\normone}[1]{||#1||_1}
\newcommand{\norminfty}[1]{||#1||_\infty}
\newcommand{\myvec}[1]{#1}

\newcommand{\Nat}{\mathbb{N}} 
\newcommand{\Zset}{\mathbb{Z}}
\newcommand{\calT}{\mathcal{T}}
\newcommand{\calM}{\mathcal{M}}

\newcommand{\ppnet}{PP-net\xspace}
\newcommand{\ppnets}{PP-nets\xspace}

\newcommand{\extract}[3]{\textsc{extract}_{#1,#2}(#3)}
\newcommand{\rstq}[2]{#1_{\downarrow #2}}
\newcommand{\incseqlen}{?}
\newcommand{\pinit}{\textsc{init}}
\newcommand{\pcov}{\textsc{cov}}
\newcommand{\pstore}{\textsc{store}}

\newcommand{\inp}{\textsc{in}}
\newcommand{\outp}{\textsc{out}}
\newcommand{\pinc}{\textsc{inc}}
\newcommand{\pdec}{\textsc{dec}}
\newcommand{\prun}{\textsc{run}}

\newcommand{\jsmall}{_\textsc{small}}
\newcommand{\jedge}{_\textsc{edge}}
\newcommand{\fmec}{_\textsc{min-cl}}
\newcommand{\fmin}{_\textsc{min-seq}}

\newcommand{\crm}[1]{\overline{#1}}
\newcommand{\restr}[1]{_{|{#1}}}

\newcommand{\setN}{\Nat}
\newcommand{\setZ}{\Zset}
\newcommand{\setQ}{\mathbb{Q}}
\newcommand{\setR}{\mathbb{R}}
\newcommand{\norm}[2]{\lVert #2 \rVert_{#1}}
\newcommand{\scalar}[2]{(#1 \cdot #2)} 


\makeatletter
\newcommand*{\da@rightarrow}{\mathchar"0\hexnumber@\symAMSa 4B }
\newcommand*{\da@leftarrow}{\mathchar"0\hexnumber@\symAMSa 4C }
\newcommand*{\xdashrightarrow}[2][]{%
  \mathrel{%
    \mathpalette{\da@xarrow{#1}{#2}{}\da@rightarrow{\,}{}}{}%
  }%
}
\newcommand{\xdashleftarrow}[2][]{%
  \mathrel{%
    \mathpalette{\da@xarrow{#1}{#2}\da@leftarrow{}{}{\,}}{}%
  }%
}
\newcommand*{\da@xarrow}[7]{%
  % #1: below
  % #2: above
  % #3: arrow left
  % #4: arrow right
  % #5: space left 
  % #6: space right
  % #7: math style 
  \sbox0{$\ifx#7\scriptstyle\scriptscriptstyle\else\scriptstyle\fi#5#1#6\m@th$}%
  \sbox2{$\ifx#7\scriptstyle\scriptscriptstyle\else\scriptstyle\fi#5#2#6\m@th$}%
  \sbox4{$#7\dabar@\m@th$}%
  \dimen@=\wd0 %
  \ifdim\wd2 >\dimen@
    \dimen@=\wd2 %   
  \fi
  \count@=2 %
  \def\da@bars{\dabar@\dabar@}%
  \@whiledim\count@\wd4<\dimen@\do{%
    \advance\count@\@ne
    \expandafter\def\expandafter\da@bars\expandafter{%
      \da@bars
      \dabar@ 
    }%
  }%  
  \mathrel{#3}%
  \mathrel{%   
    \mathop{\da@bars}\limits
    \ifx\\#1\\%
    \else
      _{\copy0}%
    \fi
    \ifx\\#2\\%
    \else
      ^{\copy2}%
    \fi
  }%   
  \mathrel{#4}%
}
\makeatother


\newcommand{\remmacro}[2]{[#1]_{#2}}
\newcommand{\rank}[1]{\operatorname{rank}(#1)}
\newcommand{\sem}[1]{\llbracket#1\rrbracket}
\newcommand{\mlcm}[1]{\operatorname{lcm}(#1)}


\newcommand*{\lastx}{}
\newcommand*{\lasty}{}
\newcommand{\drawconf}[5]{
    \xdef\lastx{#5}
    \foreach \y[count=\i] in {#3} {
        \ifnum \i = 1
            \draw[#1, #2] (\lastx,\y) -- (\lastx+#4,\y);
        \fi
        \ifnum \i > 1
            \draw[#1, #2] (\lastx,\lasty) -- (\lastx,\y);
            \draw[#1, #2] (\lastx,\y) -- (\lastx+#4,\y);
        \fi
        \xdef\lastx{\lastx+#4}
        \xdef\lasty{\y}
    }
}

\definecolor{applegreen}{rgb}{0.55, 0.71, 0.0}
\definecolor{amethyst}{rgb}{0.6, 0.4, 0.8}





%% e.g. \usepackage{tikz}
\usepackage{tikz}
\usetikzlibrary{arrows,shapes,decorations,automata,backgrounds,petri,calc,positioning}
\usepackage{amssymb,amsmath}
\usepackage{bm}
\usepackage{xspace}
\usepackage{stmaryrd}
\let\origtextsc\textsc
\renewcommand{\textsc}[1]{\textnormal{\origtextsc{#1}}}
\SetSymbolFont{stmry}{bold}{U}{stmry}{m}{n}

\usepackage{hyperref}
%%\input{myMacros.tex}
%% define non-standard environments BEYOND the ones already supplied
%% here, for example
%\theoremstyle{plain}\newtheorem{satz}[thm]{Satz} %\crefname{satz}{Satz}{S\"atze}
%% Do NOT replace the proclamation environments lready provided by
%% your own.

\def\eg{{\em e.g.}}
\def\cf{{\em cf.}}

%% due to the dependence on amsart.cls, 

%%% Theorem environments

% the following environments switch to a slanted font:
\theoremstyle{plain}

\newtheorem{thm}{Theorem}[section]
\newtheorem{cor}[thm]{Corollary}
\newtheorem{lem}[thm]{Lemma}
\newtheorem{slem}[thm]{Sublemma}
\newtheorem{prop}[thm]{Proposition}

% the following environments keep the roman font:
\theoremstyle{definition}

\newtheorem{asm}[thm]{Assumption}
\newtheorem{rem}[thm]{Remark}
\newtheorem{rems}[thm]{Remarks}
\newtheorem{exa}[thm]{Example}
\newtheorem{exas}[thm]{Examples}
\newtheorem{defi}[thm]{Definition}
\newtheorem{conv}[thm]{Convention}
\newtheorem{conj}[thm]{Conjecture}
\newtheorem{prob}[thm]{Problem}
\newtheorem{oprob}[thm]{Open Problem}
\newtheorem{oprobs}[thm]{Open Problems}
\newtheorem{algo}[thm]{Algorithm}
\newtheorem{obs}[thm]{Observation}
\newtheorem{desc}[thm]{Description}
\newtheorem{fact}[thm]{Fact}
\newtheorem{qu}[thm]{Question}
\newtheorem{oqu}[thm]{Open Question}
\newtheorem{pty}[thm]{Property}
\newtheorem{clm}[thm]{Claim}
\newtheorem{nota}[thm]{Notation}
\newtheorem{com}[thm]{Comment}
\newtheorem{coms}[thm]{Comments}

%% goal is to omit parens from note, place cite as part of theorem number
\newtheoremstyle{defC}%
  {6pt}% (space above)
  {6pt}% (space below)
  {\normalfont}% (body font)
  {}% (indent amount)
  {\bfseries}% {theorem head font}
  {{\bfseries .}}% {punctuation after theorem head}
  {5pt plus 1pt minus 1pt}% {space after theorem head}
  {\thmname{#1} \thmnumber{#2} \thmnote{\normalfont#3}}% {theorem head spec}
\theoremstyle{defC}
\newtheorem{defiC}[thm]{Definition}
\newtheorem{remC}[thm]{Remark}
\newtheorem{exaC}[thm]{Example}

\newtheoremstyle{thmC}%
  {6pt}% (space above)
  {6pt}% (space below)
  {\itshape}% (body font)
  {}% (indent amount)
  {\bfseries}% {theorem head font}
  {{\bfseries .}}% {punctuation after theorem head}
  {5pt plus 1pt minus 1pt}% {space after theorem head}
  {\thmname{#1} \thmnumber{#2} \thmnote{\normalfont#3}}% {theorem head spec}

\theoremstyle{thmC}
\newtheorem{thmC}[thm]{Theorem}
\newtheorem{propC}[thm]{Proposition}
\newtheorem{lemC}[thm]{Lemma}

% switch back to standard theorem style for document specific theorems
\AtBeginDocument{\def\shorttitle{Conservative Petri net structural liveness}}
